\chapter{State of the Art}
\label{cap:stato_arte}

This chapter reviews the perception and guidance methods relevant to
vision-based row following in agriculture. It first introduces the
operating environments and sensing alternatives, then examines how
images are converted into crop or path representations and navigation
references. Dataset preparation and learning strategies are discussed
before the integration of perception with control and experimental
validation. The final section briefly positions the thesis within this
literature, leaving implementation details to the following chapters.

\section{Autonomous navigation in agriculture: context and challenges}
\label{sec:navigazione_contesto}

\subsection{Navigation scenarios in precision agriculture}
\label{subsec:scenari_navigazione}

Autonomous agricultural navigation can be studied as a perception,
localization, planning, and guidance problem. In vision-based systems,
the camera data are processed to identify relevant structures in the
environment, such as crop rows, crop boundaries, or traversable
inter-row regions. The resulting geometric representation is then used
to derive a navigation reference and generate motion commands
\cite{bai2023vision}.

The studies considered here cover distinct agricultural settings:
greenhouse cucumber rows in Chen et al., sugar-beet fields in de Silva
et al., and vineyards and orchards in Mendez et al. and Navone et al.
\cite{chen2021navigation,silva2023rowdetection,mendez2023vineyards,
navone2025gpsfree}. Kneip et al. address the related task of detecting
the boundary between standing crop and harvested ground in wheat and
rapeseed fields \cite{kneip2020cropedge}. Following a crop row, staying
between two rows, and following a crop edge require different visual
references; they should not be treated as identical perception tasks.

\subsection{Specific challenges of row navigation}
\label{subsec:sfide_navigazione}

The reviewed studies identify difficulties associated with vegetation
structure and changing field conditions, including
\cite{silva2023rowdetection,navone2025gpsfree}:
\begin{itemize}
    \item occlusions and dense or overlapping canopies,
    \item illumination changes and shadows,
    \item presence of weeds and unstructured vegetation,
    \item missing plants, row discontinuities, and growth-stage variation.
\end{itemize}
These conditions motivate evaluating perception across different field
conditions rather than only on regular, unobstructed rows. For a
closed-loop system, that evaluation must also consider how perception
errors affect the navigation reference and the resulting robot motion.

\section{Navigation technologies: GNSS/RTK, LiDAR, and vision-based approaches}
\label{sec:localizzazione}

\subsection{Overview of navigation technologies}
\label{subsec:overview_navigation_technologies}

Agricultural navigation systems employ different sensing modalities,
including Global Navigation Satellite Systems (GNSS), Light Detection
and Ranging (LiDAR), monocular cameras, stereo
cameras, and RGB-D sensors. The choice of the sensing configuration
depends on the required localization accuracy, the structure of the
agricultural environment, and the type of navigation reference to be
extracted \cite{bai2023vision}.

\subsection{Satellite-based localization: GNSS and RTK}
\label{subsec:gnss_rtk}

In their related-work review, Navone et al. describe GNSS-based
localization as a widely adopted approach for outdoor agricultural
robots. Real-Time Kinematic (RTK) corrections support accurate
positioning for navigation and the acquisition of geo-referenced data.

However, the reliability of GNSS-based localization depends on the
quality and visibility of the received satellite signals. In agricultural
environments with thick canopies, the signal can be obstructed or
affected by multipath reflections \cite[Section~2]{navone2025gpsfree}.

GNSS/RTK provides a global positioning reference, whereas local
perception can guide motion relative to the observed crop structure.
These roles can be complementary: avoiding GNSS for row following does
not by itself solve global route planning or navigation between rows.
The following discussion focuses on local guidance when satellite
visibility is limited \cite[Sections~4.2--4.3]{bai2023vision}.

\subsection{LiDAR and stereo-vision geometry}
\label{subsec:lidar_mappe}

The related-work review of Navone et al. discusses LiDAR-based position
tracking and navigation in agricultural environments. LiDAR measurements
provide geometric information about the surroundings, including structures
such as tree trunks.

However, the suitability of LiDAR depends strongly on the structure of
the agricultural environment and on the type of sensor used.
Two-dimensional LiDAR scans can be effective in relatively simple scenarios,
such as environments where the main structures to be detected are tree
trunks. Nevertheless, 2D scans provide limited information about plants
and canopies with complex three-dimensional shapes. Multi-layer 3D LiDAR
systems can capture structures at different heights, but their additional
measurements also have to be processed by the navigation pipeline
\cite[Section~2]{navone2025gpsfree}.

Stereo vision provides another approach for obtaining geometric
information about crops and the ground. Kneip et al. use a stereo-camera
system to generate a three-dimensional point cloud and an elevation map.
The resulting height information is modeled using a Gaussian mixture,
whose parameters are estimated with Expectation--Maximization to
distinguish crop from ground. Temporal Bayesian filtering is applied
before potential crop-edge points are extracted and fitted using robust
linear regression \cite{kneip2020cropedge}.

Stereo vision and RGB-D sensing both provide geometric information, so
they are not alternatives to vision in general. For the methods reviewed
here, a useful distinction is between estimating crop geometry in a
three-dimensional or ground-plane representation and deriving a reference
directly from image-space features.

\subsection{Vision-based approaches and motivation for GPS-free methods}
\label{subsec:vision_based}

Among the sensing options described above, the following discussion
focuses on camera-based pipelines that identify crop or background
structure and derive a local guidance reference.

The motivation for GPS-free visual guidance is particularly clear in
the work of Navone et al., where dense vegetation and overlapping
canopies obstruct satellite signals and may also hide the sky at the end
of a row. Their RGB-D pipeline estimates a vegetation-based reference
without using GPS for guidance. This addresses local row following,
rather than replacing every component of a complete navigation system
\cite{navone2025gpsfree}.

The field variations introduced in
Section~\ref{subsec:sfide_navigazione} motivate the investigation of
learning-based segmentation,
but do not establish that it outperforms classical image processing in
every crop or operating condition.

The next section therefore considers how visual observations are
processed into representations suitable for extracting a guidance
reference, distinguishing classical appearance-based methods from
learned segmentation.

\section{Path perception: classical methods vs deep learning}
\label{sec:percezione_sentiero}

\subsection{Overview of path-perception methods}
\label{subsec:overview_path_perception}

Vision-based agricultural navigation methods can be distinguished
according to the image-processing strategy used to identify the relevant
crop or path structure. Classical approaches rely on color or intensity
information, thresholding, filtering, and geometric rules defined
according to the considered application. Learning-based approaches
instead use data-driven representations, including object detection and
semantic segmentation, to identify crop rows, vegetation, or traversable
regions \cite{bai2023vision}.

This distinction concerns perception, not the entire navigation pipeline:
a learned segmentation mask may still be processed by a classical
geometric algorithm. Section~\ref{sec:estrazione_traiettoria} examines
this subsequent processing stage.

\subsection{Classical color- and filter-based methods}
\label{subsec:metodi_classici_colore}

The vineyard-navigation pipeline of Mendez et al. illustrates the
combination of the following stages:
\begin{itemize}
    \item image preprocessing;
    \item segmentation through color or intensity thresholds;
    \item morphological or filtering operations to refine the result;
    \item extraction of geometric information from the segmented image.
\end{itemize}

Their pipeline applies grayscale conversion,
Gaussian filtering, Otsu thresholding, and selection of the largest blob,
followed by pixel counting to estimate the row center. The resulting
image-based reference is used to control a mobile robot in vineyard
trials \cite[Sections~3.4--4]{mendez2023vineyards}. This illustrates that a classical
pipeline can be adequate for a specific environment without requiring a
trained segmentation network.

Kneip et al. also report that classical crop-edge detection methods can
be affected by variations in illumination, shadows, weeds, barren areas,
and changes in crop appearance across growth stages
\cite[Section~2]{kneip2020cropedge}. Therefore, manually defined color- and
threshold-based methods may require suitable parameters for the
considered crop and operating conditions.

\subsection{Deep learning for semantic segmentation}
\label{subsec:deep_learning_segmentazione}

Semantic-segmentation networks assign classes to image pixels and produce
a mask used by subsequent processing. The target labels are part of the
task definition: vegetation segmentation and crop-row line prediction
produce different intermediate representations, even if both are
implemented as binary segmentation.

De Silva et al. use U-Net to predict crop-row masks in sugar-beet images.
Their annotations represent each row by a line of uniform width, rather
than labeling every plant pixel or the traversable inter-row surface.
The model is therefore trained to recover the row structure, including
gaps between plants. The authors evaluate different field conditions,
including shadows, weeds, crop sizes, and discontinuities in the rows
\cite[Sections~3.3 and~4.1]{silva2023rowdetection}.

Navone et al. apply semantic segmentation to GPS-free navigation in
dense crop rows. Their network combines a MobileNetV3 backbone with an
LR-ASPP segmentation head and produces a binary vegetation mask from
RGB images. The authors evaluate segmentation using the intersection
over union (IoU) metric, before incorporating depth information in the
navigation-reference extraction stage
\cite[Sections~3.3 and~4.1]{navone2025gpsfree}.

Compared with approaches based only on manually selected thresholds,
learning-based segmentation can represent visual patterns learned from
annotated examples. However, its performance depends on the training
data and on the similarity between the training and deployment
conditions. Variations in crop type, illumination, vegetation structure,
and field appearance therefore remain relevant challenges
\cite[Sections~2.2 and~3]{silva2023rowdetection}.


\section{From perception to a navigation reference}
\label{sec:estrazione_traiettoria}

The output of perception must be converted into a quantity usable by the
controller. In the studies reviewed here, this may be a crop-edge line,
the central crop-row line, or a horizontal image position indicating the
inter-row direction. Such references are not necessarily complete
trajectories in world coordinates. The following comparison distinguishes
the geometric operations used to obtain them.

For this comparison, it is useful to separate three questions: which
observations are selected, how their spatial arrangement is summarized,
and which quantity is passed to the controller. A segmentation mask
defines the regions or structures recognized by perception, but the
geometric stage determines how those observations become a guidance
reference. The methods below differ both in their input representation
and in this subsequent reduction: a column position, a line estimated
from candidate points, and a curve constructed from local regions retain
different amounts of spatial information.

\subsection{Global fitting and line detection}
\label{subsec:fitting_globale}
\label{subsec:hough_fitting}

Here, global fitting means estimating one model from the selected points
in an observation, not constructing a global map of the field. The
examples below contrast Hough-based path fitting with robust linear
crop-edge estimation.

Chen et al. provide a directly evaluated example using a prediction-point
Hough transform for greenhouse cucumber rows. Their pipeline selects the
lower 160 image rows, segments plants and soil using a dedicated grayscale
factor, and extracts navigation points. Least-squares regression is used
to replace a median navigation point with a predicted point, while the
angular and intersection search ranges of the Hough transform are restricted.
In the Hough parameter space, the accumulator counts intersections
between the curve associated with the predicted point and those
associated with the other navigation points, within the restricted
search range. The initial regression therefore helps constrain the
subsequent line detection; it is not simply the final navigation line.
The authors compare the result with traditional Hough and least-squares
fitting \cite[Sections~2--3]{chen2021navigation}.

Kneip et al. extract potential crop-edge locations from scanlines of an
elevation map and use robust linear regression to estimate an overall
linear crop-edge model. Each scanline is compared with a step-shaped
crop--ground transition model to locate a possible edge. Scanlines
containing only one class do not provide such a transition and are
excluded. Unlike an image-space path line, the fitted model represents
the crop boundary in ground-plane coordinates.

Global fitting provides a compact representation of the detected
structure. However, the estimated model depends on the spatial
distribution and quality of the extracted points. Outliers, missing
observations, and irregular crop structures can affect the fitted line or
model, motivating the robust estimation used by Kneip et al. Their
method uses a Huber M-estimator solved through iteratively reweighted
least squares. This separates the selection of potential edge points
from the reduction of their influence during fitting
\cite[Sections~6--7]{kneip2020cropedge}. De Silva et al. also note that Hough-based
line extraction may require parameter adjustment across crops and growth
stages \cite[Section~2.2]{silva2023rowdetection}.

\subsection{Column-wise histogram references}
\label{subsec:metodi_locali}

Histogram methods aggregate information across image columns rather
than fitting an explicit line. Mendez et al. estimate the row-center
position from the mean indices of the five columns with the largest
white-pixel counts in the selected background blob
\cite[Section~3.4.3]{mendez2023vineyards}. The meaning of the histogram
extremum therefore depends on what the mask represents.

The aggregation reduces a two-dimensional mask to a one-dimensional
profile. The resulting column position can indicate where to steer,
but does not by itself specify the slope of a path line or how the
path bends across the image. In particular, a maximum associated with
a background region and a minimum associated with vegetation are
different interpretations of the image, even when both are used to
derive a row-centering reference.

Navone et al. first combine recent vegetation masks and apply a depth
threshold to exclude distant regions. They then calculate a column-wise
sum histogram from the processed vegetation mask, rather than from a
mask of the traversable path. In SegMin, a moving average smooths the
histogram before its global minimum is located. If several columns share
the minimum, their mean position is used. This position provides a
reference for centering the robot between the crop rows.

Navone et al. also propose the SegMinD variant, which combines the
vegetation mask with normalized inverted depth so that closer elements
have a greater influence on the histogram
\cite[Section~3, including Sections~3.1--3.2]{navone2025gpsfree}.
The conversion of the histogram reference into motion commands is
discussed in Section~\ref{subsec:architetture_percezione_controllo}.

These methods provide an image-space reference without constructing a
complete field map. This does not imply that the complete pipeline is
inexpensive: sensing, neural inference where applicable, post-processing,
and control all contribute to its execution time. Their experimental
validation is discussed in Section~\ref{sec:sistemi_completi}.

\subsection{Horizontal strips and adaptive multi-ROI methods}
\label{subsec:strips}

Horizontal-strip methods extract local geometric information from
successive image regions before constructing a navigation reference.
Li et al. provide a directly relevant example: their Faster-U-net model
predicts the inter-row path, and an adaptive multi-ROI procedure extracts
central points from the predicted binary mask. Horizontal bands define
incremental sub-ROI search regions; the resulting points are used as
control points for a multisegment cubic B-spline navigation curve.

Li et al. also discuss how the height of the search regions changes the
number of available points. Larger search regions produce fewer
control points; smaller regions produce more points, with a risk of
overfitting discussed by the authors. They select a height of 40 pixels
in their experimental configuration, illustrating that strip size and
fitting must be considered together rather than as independent choices
\cite[Section~2.2.3, Figure~9]{li2022multiroi}.

Unlike a single histogram over the processed image, strip-wise
extraction preserves the vertical ordering of local observations: each
point belongs to a different part of the image. The sequence can
therefore represent changes in the observed path position before a
geometric model is fitted. Subdividing the image alone, however, does
not determine which region to select when several candidates are
present. The interpretation of the mask and the rule for obtaining a
representative point remain part of the method.

This approach distinguishes local point extraction from the subsequent
geometric model: using horizontal strips does not require fitting a
straight line, just as fitting a single line does not rule out extracting
its input points locally. Strip dimensions, candidate selection, and
model choice must be evaluated within the complete pipeline; the use of
strips alone does not establish robustness or computational efficiency.

\subsection{Triangle-scan extraction of the central crop row}
\label{subsec:triangle_scan}

The method proposed by de Silva et al. is the Triangle Scan Method (TSM).
It first searches a horizontal strip at the top of the predicted crop-row
mask. The column with the largest normalized pixel sum defines an upper
anchor point; a preset anchor is used when the scan fails a validation
threshold. Candidate lines are then scanned from this anchor toward
positions on the base of a triangular region of interest. The line with
the highest mask-pixel sum determines the lower endpoint of the central
crop row \cite[Section~4.2]{silva2023rowdetection}.

The two endpoints provide an angular error and an image-space position
error for visual servoing. This is different from independently extracting
a point in each of several horizontal strips and subsequently fitting a
line through those points. It also follows a crop-row representation,
not a mask of traversable ground.

Thus, the horizontal strip used by TSM has a different role from the
successive regions used by a multi-ROI fitting method. It establishes
one endpoint from which candidate lines are evaluated; it does not
produce a sequence of independent control points along the path.
The preset anchor provides a fallback when the upper scan is not
accepted, while the triangular ROI bounds the subsequent search.
These choices constrain the family of lines that can be selected
\cite[Section~4.2, Algorithm~1]{silva2023rowdetection}.

\subsection{Comparison of representations and geometric assumptions}
\label{subsec:confronto_riferimenti}

Table~\ref{tab:confronto_riferimenti_geometrici} summarizes the information
flow in the reviewed methods. The comparison separates the representation
provided by perception from the geometric operation and its output.
It is not a ranking: following a crop edge, identifying a central crop
row, and estimating an inter-row path are different tasks, even when
their outputs are subsequently used for steering.

\begin{table}[p]
    \centering
    \footnotesize
    \setlength{\tabcolsep}{3pt}
    \renewcommand{\arraystretch}{1.2}
    \caption{Perceptual inputs and geometric navigation references in the
    reviewed studies. RGB-D denotes color and depth sensing. The table
    compares representations, not performance under a common benchmark.}
    \label{tab:confronto_riferimenti_geometrici}
    \begin{tabularx}{\textwidth}{@{}>{\raggedright\arraybackslash}p{0.14\textwidth}
        >{\raggedright\arraybackslash}p{0.10\textwidth}
        >{\raggedright\arraybackslash}X
        >{\raggedright\arraybackslash}X
        >{\raggedright\arraybackslash}p{0.18\textwidth}@{}}
        \toprule
        Study & Sensor & Perceptual input & Geometric operation & Output reference \\
        \midrule
        Chen et al.\ \cite{chen2021navigation}
        & RGB
        & Navigation points from plant--soil segmentation
        & Regression-based predicted point and restricted Hough voting
        & Image-space path line \\
        \addlinespace
        Kneip et al.\ \cite{kneip2020cropedge}
        & Stereo
        & Crop--ground classes in an elevation map
        & Scanline edge candidates and robust linear regression
        & Crop-edge line in ground-plane coordinates \\
        \addlinespace
        Mendez et al.\ \cite{mendez2023vineyards}
        & RGB
        & Largest background blob after thresholding
        & Column counts and mean indices of the five largest counts
        & Horizontal row-center position \\
        \addlinespace
        Navone et al.\ \cite{navone2025gpsfree}
        & RGB-D
        & Temporally combined vegetation masks refined with depth
        & Smoothed column histogram; SegMinD also weights by inverted depth
        & Position of the histogram minimum \\
        \addlinespace
        Li et al.\ \cite{li2022multiroi}
        & RGB
        & Predicted inter-row path mask
        & Local points from horizontal sub-ROIs and cubic B-spline construction
        & Curved image-space path reference \\
        \addlinespace
        De Silva et al.\ \cite{silva2023rowdetection}
        & RGB
        & Predicted masks of crop-row lines
        & Upper anchor search followed by triangular line scanning
        & Central row line defined by two endpoints \\
        \bottomrule
    \end{tabularx}
\end{table}
\clearpage

A central modeling choice is whether the selected observations can be
represented adequately by a straight line over the region being
processed. Kneip et al. explicitly assume a predominantly linear crop
edge \cite[Section~7]{kneip2020cropedge}, while Li et al. motivate the
use of a B-spline through curved crop rows for which straight-line
fitting can give larger errors \cite[Section~2.2.3]{li2022multiroi}.
Linearity must be interpreted in the coordinate system of the method:
an image-space path line and a crop-edge line in the ground plane do
not represent the same measured quantity.

The distinction also separates outlier handling from model choice.
Reducing the influence of inconsistent points can help estimate a line
from noisy observations, but does not turn that line into a model of
curvature. Conversely, a curved representation can describe changes in
direction, but its suitability still depends on the selected points
and on whether their variation reflects the path rather than errors
in perception. This is a methodological consequence of the compared
representations, not evidence that one family is universally more
accurate or less expensive.

Finally, local point extraction and global fitting are not mutually
exclusive. Points obtained separately from strips or scanlines can
still be combined into one model for the current observation. For a
navigation system, the relevant choice is therefore not simply
``local versus global'': it is the combination of segmentation target,
selection rule, geometric model, coordinate system, and control input.
The experimental evidence for these combinations is considered in
Section~\ref{subsec:validation_scope}, separately from their geometric
description.


\section{Datasets, transfer learning, and annotation}
\label{sec:dataset_transfer}

The segmentation-based methods discussed above depend on the data used
to learn their intermediate representations. This section considers the
availability of annotations, the initialization of network weights, and
model-assisted labeling as distinct aspects of dataset and model preparation.

\subsection{Annotation challenges in agricultural datasets}
\label{subsec:scarsita_dataset}

Supervised training of semantic-segmentation networks uses images associated
with pixel-level target masks. In agricultural applications, obtaining
these annotations can be costly in terms of time and human effort
\cite{navone2025gpsfree}.

This issue is particularly relevant for crop-row navigation, since the
visual appearance of agricultural environments can vary according to the
crop type, plant structure, illumination, and field conditions. A dataset
collected in one environment may therefore not represent all the
conditions encountered during robot operation. The dataset of de Silva
et al., for example, explicitly represents 11 field variations, including
growth stages, light conditions, weed densities, and discontinuities
\cite[Section~3]{silva2023rowdetection}.

\subsection{Transfer learning for agricultural segmentation}
\label{subsec:transfer_learning}

The form of transfer learning considered here initializes a network with
weights learned on a previous dataset and adapts it to agricultural
segmentation. Navone et al. start from ImageNet-pretrained weights and
train crop-specific models using synthetic AgriSeg images supplemented
by 100 real images of miscellaneous crops, different from the real test
sets \cite[Section~4.1, Table~1]{navone2025gpsfree}.
Initializing from pretrained weights and combining synthetic with real
training images are separate choices. The reported training configuration
should not be interpreted as isolating their individual effects.

\subsection{SAM 3 and model-assisted annotation}
\label{subsec:auto_labeling}

Model-assisted annotation is distinct from transfer learning: the former
concerns generating training labels, whereas the latter concerns adapting
model weights.

Segment Anything Model~3 (SAM~3), developed by Meta, supports
promptable concept segmentation. Short text phrases, image exemplars, or their
combination specify a concept, and the model predicts masks for matching
instances in images or videos. To build the model's training data, the
authors use a data engine that combines mask proposals with checks of
quality and completeness by human and automated verifiers
\cite[Sections~2 and~4]{carion2026sam3}. This data-collection process is
distinct from simply running the released model on a new image.

For agricultural dataset preparation, SAM~3 offers a possible alternative
to manually labeling every image. However, the authors report limitations
in recognizing fine-grained concepts outside the training domain and in
interpreting ambiguous text prompts
\cite[Section~2 and Appendix~C]{carion2026sam3}.

This thesis investigates whether SAM~3 can generate plant and soil masks
of sufficient quality to serve as training labels for a segmentation
model, while reducing the time required to prepare the dataset. The aim
is to assess how much manual labeling can be avoided in the agricultural
images considered here. The evaluation therefore needs to consider both
mask quality and the time spent generating, checking, and, where necessary,
correcting the labels. Faster dataset preparation is an intended benefit
to be evaluated, rather than an assumed outcome.

\section{Complete systems: perception, control, and validation}
\label{sec:sistemi_completi}

\subsection{Perception--control architectures in agricultural robots}
\label{subsec:architetture_percezione_controllo}

Complete agricultural navigation systems combine environmental
perception, geometric processing, and robot control. A typical pipeline
consists of acquiring images or depth data, extracting crop-row or
traversable-path information, deriving a local navigation reference, and
converting the reference into motion commands.

Navone et al. implement this type of architecture using an RGB-D camera.
The vegetation-mask and histogram processing described in
Section~\ref{subsec:metodi_locali} provides the navigation reference.
Its displacement from the image center is used to compute the linear
and angular velocity commands through custom functions. An
Exponential Moving Average is subsequently applied to smooth the velocity
commands and reduce abrupt changes in the robot motion
\cite[Section~3.4]{navone2025gpsfree}.

Other studies use different control inputs. Mendez et al. switch between
two proportional control modes based on row-center displacement and the
distribution of background pixels on either side of the detected center,
while keeping linear velocity constant
\cite[Section~3.5]{mendez2023vineyards}. De Silva et al. use a weighted
combination of the central row's angular and positional errors in a
proportional visual-servoing controller, with constant forward velocity,
for their simulation experiments. This is distinct from the
image-Jacobian-based IBVS controller discussed as a comparison in their
paper \cite[Section~4.3, Eq.~(5)]{silva2023rowdetection}. These examples show
why the representation extracted by perception and the quantities used
by control must be considered together.

\subsection{Scope of the experimental evidence}
\label{subsec:validation_scope}

The reviewed papers do not all validate the same level of the navigation
pipeline. Kneip et al. evaluate crop-edge detection on recorded harvesting
data, including comparison with manually labeled edge lines; this is
evidence about crop-edge estimation, rather than inter-row tracking
\cite[Section~8]{kneip2020cropedge}. Chen et al. compare image-based
path-fitting methods, including their angular deviations and processing
times. Their timing comparison covers prediction-point Hough, traditional
Hough, and least-squares fitting. It concerns the path-fitting stage in
the reported experiment, not the latency of a complete autonomous
navigation system \cite[Section~3.6, Tables~6--7]{chen2021navigation}.

De Silva et al. evaluate row detection on real images and visual servoing
in simulation. Their detection evaluation concerns image-space line
estimation, whereas the simulated visual-servoing experiments evaluate
the resulting robot motion. They
also mention preliminary real-world row following over 6~m, while
identifying further field validation as future work
\cite[Sections~5.2--5.3 and~6, Table~5]{silva2023rowdetection}.

Mendez et al. report closed-loop navigation trials in a real vineyard,
with LiDAR used to measure performance rather than to generate the
camera-based navigation reference. Their plotted centering error is the
difference between the measured distances on the two sides of the row;
it should not be read as an identically defined metric to every other
study's cross-track error \cite[Sections~3.6 and~4]{mendez2023vineyards}.

Navone et al. evaluate SegMin and SegMinD both in simulation and in
real crop fields. For the field trials, they estimate lateral
displacement using LiDAR and assess heading using an attitude sensor
and a row-direction reference from satellite images; a ground-truth
trajectory is not available. Thus, GPS-free visual guidance should not
be confused with camera-only measurement of experimental performance
\cite[Sections~4.3--4.4]{navone2025gpsfree}.

Consequently, segmentation quality, image-space line accuracy, metric
tracking error, and processing time should be reported separately.
Comparisons across studies require attention to the sensors, datasets,
hardware, and experimental conditions, rather than a ranking based on
unlike metrics.

\section{Final considerations and positioning of the thesis}
\label{sec:considerazioni_finali_stato_arte}

The reviewed methods illustrate different ways of connecting perception
to guidance. Column-wise histograms provide a horizontal reference
without explicitly fitting a path line
(Section~\ref{subsec:metodi_locali}), whereas triangle scanning identifies
a central crop-row line from a row-mask representation
(Section~\ref{subsec:triangle_scan}). Crop-edge fitting instead estimates
the boundary between standing crop and harvested ground. These
representations address different tasks and are not interchangeable
solely because each can supply a controller with a guidance reference.

Among the reviewed approaches, the adaptive multi-ROI method of Li
et al.\ \cite{li2022multiroi} provides the closest architectural comparison
for the geometric stage investigated in this thesis: both extract local
points from a segmentation mask using horizontal regions before
constructing a navigation reference. However, Li et al. segment the
inter-row path and use the extracted points to construct a multisegment
cubic B-spline. This similarity concerns the processing structure,
not equivalence of the algorithms or their experimental performance.

The thesis investigates semantic segmentation of plants and soil,
followed by selection of candidate soil intervals in successive strips.
Local class support and spatial continuity guide interval selection;
the selected interval midpoints then contribute to a weighted
straight-line fit with residual-based reweighting. The resulting
reference is defined in image coordinates and adopts a locally straight
path model, rather than the curved representation used by Li et al.
Segmenting visible soil does not itself identify a unique navigation
corridor: corridor selection belongs to the geometric stage, and the
resulting line does not constitute a guarantee of obstacle-free motion.

This positions the work as the design and evaluation of a pipeline that
integrates segmentation, geometry, and control, rather than the
introduction of horizontal strips as a new general technique. Dataset preparation,
including manual and model-assisted annotation, and the choices made
in point selection, fitting, and control are developed in the following
chapters. Their evaluation must distinguish annotation quality,
segmentation accuracy, geometric-reference quality, and closed-loop
navigation performance; advantages in robustness or computational cost
cannot be inferred from the architecture alone.